\documentclass[11pt,a4paper]{article}
\usepackage[T1]{fontenc}
\usepackage{mathpazo,microtype,mathtools,amssymb,amsthm,booktabs,array}
\usepackage[margin=25mm]{geometry}
\usepackage{xurl,hyperref,listings}
\hypersetup{hidelinks,pdftitle={An explicit degree-29984 Smale counterexample}}
\renewcommand{\UrlFont}{\normalfont\ttfamily\small}
\lstset{basicstyle=\ttfamily\small,breaklines=true,columns=fullflexible,frame=single}
\newtheorem{theorem}{Theorem}\newtheorem{lemma}{Lemma}\newtheorem{proposition}{Proposition}
\newcommand{\C}{\mathbb C}\newcommand{\Q}{\mathbb Q}\newcommand{\Smal}{\mathcal S}
\newcommand{\abs}[1]{\left|#1\right|}
\newcommand{\file}[1]{\textnormal{\small\ttfamily\nolinkurl{#1}}}
\newcommand{\astar}{\widehat a}
\allowdisplaybreaks[2]\emergencystretch=2em
% Generated from the completed, audited interval certificates.
\newcommand{\CertifiedLower}{1.00000049595801}
\newcommand{\CertifiedUpper}{1.00000049595804}
\newcommand{\BaseLower}{1.00000000407528}
\newcommand{\BaseUpper}{1.00000052222008}
\newcommand{\RefinedLower}{1.00000049595801828345}
\newcommand{\RefinedUpper}{1.00000049595803927317}
\newcommand{\RefinedCount}{316}
\newcommand{\MinimizerIndex}{28891}
\newcommand{\OtherLower}{1.0000004992776471}
\newcommand{\UniquenessStatement}{Exactly one critical point attains the minimum.}
\newcommand{\UniquenessExplanation}{This lower bound exceeds the upper bound for the minimizing disk, so the minimizer is unique.}

\title{An explicit polynomial of degree 29,984\\in Smale's mean value problem}
\author{}
\date{}
\begin{document}\maketitle
\begin{abstract}
We specify a Gaussian-rational polynomial of degree 29,984 whose normalized
divided difference exceeds one at every critical point. Its normalization
point is independent of the anchor of the analytic model. A complete interval
calculation isolates all 29,983 critical points and encloses their normalized
quotients. We give the finite definition, local root and value estimates,
and outward-rounding arguments used by the verifier. A local comparison
shows that the same polynomial has a critical quotient below one at its model
anchor.
\end{abstract}

\section{Result and scope}
\begin{theorem}\label{thm:result}
The exact polynomial $p\in\Q(i)[x]$ defined in Section~\ref{sec:definition},
with the finite coefficient tables in \file{polynomial.json}, satisfies
\[
 \deg p=29\,984,\qquad p(0)=0,\qquad p'(0)=1,
\]
and
\begin{equation}\label{eq:result}
 \boxed{\CertifiedLower<\Smal(p,0):=
 \min_{p'(c)=0}\abs{\frac{p(c)}c}<\CertifiedUpper.}
\end{equation}
All 29,983 critical points are simple. \UniquenessStatement
\end{theorem}

The finite specification and the rational containing disks suffice to reproduce
the interval proof. The computation isolates and checks every critical point
of this polynomial.

The model anchor $a$ and the normalization base $\astar$ are different.
The letter $b$ denotes the rational approximation to an algebraic constant.
Every numerical string, including terminating decimals and scientific
notation, denotes an exact rational number. The specified phase has modulus
one exactly.

\section{Explicit finite polynomial}\label{sec:definition}
Set
\[
 m=16,\quad N=1874,\quad L=N-2=1872,\qquad
 A=\frac{1061603}{250000},\quad C=1+A^2,\quad S=A+1/A,
\]
\[
 q(z)=\frac{z^{16}+iA}{S},\qquad v(z)=iq(z).
\]
The exact amplitude, Cayley parameter, and phase are
\[
 \gamma=0.61480744530374742407,\qquad t=55.5578744900222005754,\qquad
 \omega=\frac{1+it}{1-it}.
\]
Thus $\gamma>0$ and $|\omega|=1$ exactly. The normalization base is
\[
 \boxed{\astar=1.04015231339331770677+0.302102329270430736141\,i.}
\]
The other Gaussian-rational parameters $a,r,b,\xi_1,\xi_2$ are printed in the
appendix and supplied in the JSON specification.

Define
\[
 d_0=C+\bar r^{16}(-A^2+iA),\quad
 d_1=iA-\bar r^{16}(1+iA),\quad
 \alpha=-\frac{Cd_1}{iAd_0}.
\]
For $1\le k\le15$, let $\nu_k=1-k/16$ and define
\[
 b_{k,0}=1,\quad b_{k,j}=b_{k,j-1}\frac{\nu_k-j+1}{j},\qquad
 c_{k,0}=1,\quad c_{k,j}=\alpha c_{k,j-1}+b_{k,j}\quad(1\le j\le L).
\]
Thus $b_{k,j}=\binom{\nu_k}{j}$. With
\[
 J(z)=-rC(1+v(z))+\bar r^{15}(1+iA)z^{16},
\]
define
\begin{equation}\label{eq:h}
\begin{split}
 h_L(z)={}&\frac{J(z)}{d_0}\sum_{j=0}^{L-1}(\alpha v(z))^j\\
 &+\frac{(1-|r|^2)C}{d_0}\sum_{k=1}^{15}\bar r^{k-1}b^kz^k
                   \sum_{j=0}^{L}c_{k,j}v(z)^j.
\end{split}
\end{equation}
The first sum ends at $L-1$, the other sums at $L$. The $b$ in this formula
is the specified rational complex number, not an unspecified algebraic root.

For the corrections, set
\[
 K(z)=(z-\xi_1)^2(z-\xi_2)^2,\qquad \mu=(\xi_1+\xi_2)/2,
\]
\[
 R_0(z)=\sum_{\ell=0}^{63}\sum_{k=0}^{15}
 d_{\ell k}\left(\frac z{6/5}\right)^kq(z)^\ell,
 \qquad T(z)=\sum_{j=0}^3\eta_j(z-\mu)^j,
\]
\[
 \Delta_0(z)=(z-a)^2[K(z)R_0(z)+T(z)],\qquad
 X(z)=(z-a)^2\sum_{\ell=0}^{127}\sum_{k=0}^{15}
 x_{\ell k}\left(\frac z{6/5}\right)^kq(z)^\ell.
\]
The arrays $d,\eta,x$ are respectively \file{correction_coefficients},
\file{free_coefficients}, and \file{extra_correction_coefficients} in the
specification. The retained Hermite coefficient is exactly zero.
The high-degree tail is
\[
 E(z)=\sum_{\ell=0}^{7}\sum_{k=0}^{15}
 t_{\ell k}\left(\frac z{6/5}\right)^kq(z)^\ell,
 \qquad H(z)=\omega[\gamma q(z)^{1874}+q(z)^{1866}E(z)],
\]
where $t_{\ell k}$ is the table \file{tail_coefficients}, not the scalar Cayley
parameter $t$. Define
\begin{equation}\label{eq:P}
 P(z)=(z-a)\frac{h_L(z)}{h_L(a)}+\Delta_0(z)+X(z)+H(z)
\end{equation}
and finally
\begin{equation}\label{eq:p}
 \boxed{p(x)=\frac{P(\astar+x)-P(\astar)}{P'(\astar)}.}
\end{equation}
Both corrections vanish to order two at $a$, whereas normalization is at
$\astar$. The evaluator uses the full finite values $P(\astar)$ and
$P'(\astar)$ in \eqref{eq:p}, and encloses the latter away from zero.
This gives $p(0)=0$ and $p'(0)=1$ exactly.

The degree bounds are
\[
\begin{array}{c|r}
\text{term}&\text{degree, at most}\\\hline
 h_L&29\,967\\(z-a)h_L&29\,968\\
 \Delta_0&1029\\X&2049\\q^{1866}E&29\,983.
\end{array}
\]
Only $\omega\gamma q^{1874}$ contributes in degree 29,984. Its coefficient
is nonzero, so translation and the nonzero normalization in \eqref{eq:p}
preserve degree exactly.

\section{An independent base point}\label{sec:search}
The anchor $a$ defines the analytic model, while the critical-point functional
can be evaluated at any base point $b_0$ with $P'(b_0)\ne0$. For a fixed
polynomial and fixed critical point $c$, define
\[
 Q_c(b_0)=\frac{P(c)-P(b_0)}{(c-b_0)P'(b_0)}.
\]
Where the denominators and numerator are nonzero,
\begin{equation}\label{eq:base-gradient}
 \frac{d}{db_0}\log Q_c(b_0)
 =-\frac{P'(b_0)}{P(c)-P(b_0)}+\frac1{c-b_0}
  -\frac{P''(b_0)}{P'(b_0)}.
\end{equation}
Changing $b_0$ leaves every critical point fixed. The base point therefore
supplies two real parameters without changing the degree. The normalization
in \eqref{eq:p} follows after the polynomial and base point are chosen.





The local computation in \file{base_comparison.log} isolates a root of $P'$
in containing disk 27158 for which
\[
 0.99982426177940
 <\abs{\frac{P(c)-P(a)}{(c-a)P'(a)}}
 <0.99982426177944<1.
\]
Thus the independent base point is essential for this example: at the model
anchor $a$, the same polynomial has a critical quotient below one.
\section{Complete critical-point coverage}\label{sec:count}
Computations use the $P$-coordinate. A critical point $c$ corresponds to
$c-\astar$ for $p$, and the required quotient is
\begin{equation}\label{eq:quotient}
 Q(c)=\abs{\frac{P(c)-P(\astar)}{(c-\astar)P'(\astar)}}.
\end{equation}
The file \file{critical_disks.txt} supplies exactly $29\,983$ rational centers
and positive rational radii. These are \emph{containing disks}; the accepted
root disks may be smaller and may have recentered dyadic centers.

\begin{lemma}[Local linear comparison]\label{lem:model}
Let $G,g$ be holomorphic on a neighborhood of $D=\{z:|z-z_0|\le\rho\}$.
Suppose $E\ge\sup_D|G-g|$ and $V\ge\sup_D|g'(z)-g'(z_0)|$. If
\begin{equation}\label{eq:rouche}
 |g(z_0)|+\rho V+E<\rho|g'(z_0)|,
\end{equation}
then $G$ has exactly one zero in the open disk, counted with multiplicity.
\end{lemma}
\begin{proof}
On the boundary, integration along the radial segment gives
\[
 |G(z)-g'(z_0)(z-z_0)|\le E+|g(z_0)|+\rho V
                  <|g'(z_0)(z-z_0)|.
\]
Rouch\'e's theorem compares $G$ with a nonzero linear polynomial, which has
one zero inside~\cite{DLMFComplex}. In particular, the unique zero is simple.
\end{proof}

The first local test applies this lemma with $G=P'$ and with $g$ the derivative
of the analytic comparison function defined below. When its approximation
error forces a needlessly large disk, the code uses finite-polynomial values
at the center instead.

\begin{lemma}[Finite-center comparison]\label{lem:finite}
Let $e=P'-g$ be holomorphic on a neighborhood of the disk of radius $R>\rho$
about $z_0$, and suppose $|e|\le E_R$ there. If
$V\ge\sup_{|z-z_0|\le\rho}|g'(z)-g'(z_0)|$, then
\[
 V_P=V+\frac{2\rho E_R}{(R-\rho)^2}
 \ \ge\ \sup_{|z-z_0|\le\rho}|P''(z)-P''(z_0)|.
\]
Consequently,
\begin{equation}\label{eq:finite}
 |P'(z_0)|+\rho V_P<\rho|P''(z_0)|
\end{equation}
proves exactly one simple critical point in the disk.
\end{lemma}
\begin{proof}
Cauchy's derivative estimate yields $|e''(z)|\le2E_R/(R-\rho)^2$ on the
smaller disk. Integrate along a radius to bound $e'(z)-e'(z_0)$, add the
bound for $g'$, and apply the same linear comparison to $P'$.
\end{proof}

Every containing-disk pair that is not already separated by its real-coordinate
gap is checked using a rigorous lower bound for center distance and an upper
bound for the sum of radii. The real-coordinate sweep uses the largest radius,
so omitted pairs are provably disjoint. Every accepted root disk is checked
to lie inside its corresponding containing disk. Finite Newton steps merely
propose centers: only the subsequent strict inequalities certify a root.

There are $29\,983$ disjoint accepted disks, each containing one root with
multiplicity, and $\deg P'=29\,983$. The fundamental theorem of algebra
therefore exhausts the critical set and proves that every critical point
is simple. The sixteen work partitions are index classes modulo sixteen.

\section{Analytic comparison at the fixed model anchor}\label{sec:model}
Set $D=\Delta_0+X$ and introduce, solely for the proof,
\[
 \beta=\left(\frac{1+iA}{C}\right)^{1/16},\quad U=1+iq,\quad
 W=\beta zU^{-1/16},\quad B(w)=\frac{w-r}{1-\bar r w},
\]
\[
 B_a=B(W(a)),\qquad G(z)=\frac{B(W(z))}{B_a},\qquad
 P_*(z)=(z-a)G(z)+D(z)+H(z),\qquad g=P_*'.
\]
Fractional powers have their principal values. In every comparison box, the
code excludes the branch cut and all relevant poles. The segment-denominator
check in Section~\ref{sec:remainder} permits these local comparisons also when
$|q|\ge1$. The disjoint-disk degree count supplies complete critical-point
coverage.
Direct differentiation gives
\[
 W'=\frac{\beta}{C}U^{-1-1/16},\quad
 G'=\frac{(1-|r|^2)W'}{(1-\bar rW)^2B_a},\quad
 G''=G'\left[-\frac{17iAz^{15}}{1+iAz^{16}}
                    +\frac{2\bar rW'}{1-\bar rW}\right].
\]
The exact rational identity $W^{16}=(1+iA)z^{16}/(1+iAz^{16})$ gives
\begin{equation}\label{eq:rationalization}
 B(W)=\frac{J}{d_0(1-\alpha v)}+
 \frac{(1-|r|^2)C}{d_0}\sum_{k=1}^{15}\bar r^{k-1}\beta^kz^k
                       \frac{(1+v)^{\nu_k}}{1-\alpha v}.
\end{equation}
This is obtained by multiplying $B(W)$ by
$\sum_{j=0}^{15}(\bar rW)^j$. It is finite algebra, not a convergence assertion
about that geometric series.

\subsection{The algebraic constant is rigorously isolated}
Let $\epsilon=2\cdot10^{-60}$ and $T_\beta(z)=z^{16}-(1+iA)/C$.
Using exact GMP rational arithmetic, the initializer verifies
\begin{multline*}
 |\Re T_\beta(b)|+|\Im T_\beta(b)|+
 \sum_{j=2}^{16}\binom{16}{j}(|\Re b|+|\Im b|)^{16-j}\epsilon^j\\
 <16(\Re b)^{15}\epsilon.
\end{multline*}
Rouch\'e's theorem therefore places one root within $\epsilon$ of $b$.
The verified inequalities $\Re b>\epsilon$, $\Im b>\epsilon$, and
$10(\Im b+\epsilon)<\Re b-\epsilon$ select the principal root. In particular,
\[
 |b^k-\beta^k|\le a_k:=k\epsilon(|\beta|+\epsilon)^{k-1}.
\]
The polynomial itself still uses $b\in\Q(i)$ throughout.

\section{Finite binomial remainders beyond the Taylor circle}\label{sec:remainder}
For $0<\nu<1$, define
\[
 T_L^{\nu}(v)=(1+v)^\nu-\sum_{j=0}^L\binom\nu j v^j.
\]
\begin{lemma}\label{lem:stieltjes}
For $v\notin(-\infty,-1]$,
\begin{equation}\label{eq:stieltjes}
 T_L^{\nu}(v)=\frac{(-1)^L\sin(\pi\nu)}{\pi}\,v^{L+1}
             \int_0^1\frac{t^{L-\nu}(1-t)^\nu}{1+tv}\,dt.
\end{equation}
Set $u=|v|$, $\sigma(v)=\min_{0\le t\le1}|1+tv|$, and
$b_L=|\binom\nu{L+1}|$. Then
\begin{align}
 |T_L^{\nu}(v)|&\le b_Lu^{L+1}/\sigma,\label{eq:T0}\\
 |(T_L^{\nu})'(v)|&\le b_L\left(
 \frac{(L+1)u^L}{\sigma}
 +\frac{L+1-\nu}{L+2}\frac{u^{L+1}}{\sigma^2}\right).
 \label{eq:T1}
\end{align}
These formulas do not require $|v|<1$.
\end{lemma}
\begin{proof}
For $|v|<1$, expand $(1+tv)^{-1}$ and apply Euler's beta integral
\cite{DLMFBeta}. Its coefficients are exactly the binomial tail coefficients,
including their signs. Both sides of \eqref{eq:stieltjes} are holomorphic on
the slit plane, so analytic continuation gives the identity there. The
integral's total mass, after multiplication by $\sin(\pi\nu)/\pi$, is $b_L$.
Its normalized first moment in $t$ is $(L+1-\nu)/(L+2)$. Taking absolute
values, and then differentiating, proves \eqref{eq:T0}--\eqref{eq:T1}.
\end{proof}

For a complex box with point midpoint $c$ and enclosing disk radius $r_0$,
$\sigma(v)\ge\sigma(c)-r_0$. The value $\sigma(c)$ is the distance from zero
to the segment from $1$ to $1+c$: it is $1$, $|1+c|$, or
$|\Im c|/|c|$, according as the minimizing projection lies at an endpoint or
in the interior. When interval tests cannot determine the case uniquely,
the smaller candidate bound is used. A strict positive lower bound both
bounds the integral and excludes the binomial branch cut on the entire box.

The recurrence for $c_{k,j}$ implies the finite identity
\[
 \mathcal R_{k,L}(v):=\frac{(1+v)^{\nu_k}}{1-\alpha v}
                  -\sum_{j=0}^Lc_{k,j}v^j
 =\frac{T_L^{\nu_k}(v)+\alpha c_{k,L}v^{L+1}}{1-\alpha v}.
\]
With $d=|1-\alpha v|$, take
\[
 T_{k,0}=\frac{b_Lu^{L+1}/\sigma+|\alpha c_{k,L}|u^{L+1}}d,
\]
\[
 T_{k,1}=\frac{b_L[(L+1)u^L/\sigma+
             (L+1-\nu_k)u^{L+1}/((L+2)\sigma^2)]
              +(L+1)|\alpha c_{k,L}|u^L}d
               +\frac{|\alpha|T_{k,0}}d.
\]
They bound $|\mathcal R_{k,L}|$ and $|\mathcal R_{k,L}'|$ respectively.
All denominator lower bounds are checked to be positive.

\subsection{Combining all analytic errors}
Set $B_0=|1+v|$, $t=|\alpha|u$, and
\[
 w_k=\frac{(1-|r|^2)C|r|^{k-1}}{|d_0|},\qquad
 H_k=B_0^{\nu_k}/d,\qquad
 H_{k,1}=\nu_kB_0^{\nu_k-1}/d+|\alpha|B_0^{\nu_k}/d^2.
\]
The finite geometric identity and \eqref{eq:rationalization} give
\begin{align*}
 e_h={}&\frac{|J|t^L}{|d_0|d}+
 \sum_{k=1}^{15}w_k|z|^k(|b|^kT_{k,0}+a_kH_k),\\
 e_{h'}={}&\frac{|J'|t^L}{|d_0|d}+
 \frac{|J||\alpha||v'|}{|d_0|}
             \left(\frac{Lt^{L-1}}d+\frac{t^L}{d^2}\right)\\
 &+\sum_{k=1}^{15}w_k\left[
 k|z|^{k-1}(|b|^kT_{k,0}+a_kH_k)
 +|z|^k|v'|(|b|^kT_{k,1}+a_kH_{k,1})\right].
\end{align*}
These bound $|h_L-B(W)|$ and the difference of their derivatives.
The initializer proves $|h_L(a)-B_a|<\epsilon_a=10^{-50}$ and
$h_a^-:=|B_a|-\epsilon_a>0$. Hence, for $G_L=h_L/h_L(a)$,
\[
 |G_L-G|\le e_g:=\frac{e_h+|G|\epsilon_a}{h_a^-},\qquad
 |G_L'-G'|\le e_{g'}:=\frac{e_{h'}+|G'|\epsilon_a}{h_a^-}.
\]
It follows that $|P'-g|\le e_g+|z-a|e_{g'}$. The polynomial corrections and
high-degree tail cancel exactly from this difference. Every expression is
evaluated on the entire current box, not on a set of samples.

\section{Finite-polynomial evaluation at a center}\label{sec:evaluation}
The code has two implementations of $h_L,h_L',h_L''$. Direct finite Horner
evaluation propagates values and two derivatives through the rational
coefficient recurrence. The faster implementation subtracts a rigorously
enclosed \emph{finite} remainder from $(1+v)^\nu$.

Normalize the beta weight in \eqref{eq:stieltjes} to a probability measure and
define $\mu_j=\mathbb E[(1-t)^j]=(\nu+1)_j/(L+2)_j$. A finite geometric identity
implies, for any integer $K\ge0$,
\[
 \mathbb E\frac1{1+tv}=\sum_{j=0}^{K-1}
 \frac{\mu_jv^j}{(1+v)^{j+1}}+\varepsilon_K,
 \qquad
 |\varepsilon_K|\le
 \frac1{\sigma(v)}\left(\frac{|v|}{|1+v|}\right)^K\mu_K.
\]
For $s=1-t$, the finite geometric remainder is
$[sv/(1+v)]^K/(1+tv)$, which proves the bound directly.

The implementation selects the best bound among $0\le K\le32$ on a slightly
larger box. If $E$ bounds the error in $T_L^\nu$ there, and a disk of radius
$R$ about every center-box point lies in that larger box, Cauchy's estimate
bounds the errors in its first two derivatives by $E/R$ and $2E/R^2$.
The branch and denominator checks are repeated on the larger box.

For clarity, let $S_K(v)$ denote the finite moment sum and write
$A_{k,L}(v)=\sum_{j=0}^Lc_{k,j}v^j$. With $e_k=\alpha c_{k,L}$,
\[
 A_{k,L}=\frac{(1+v)^{\nu_k}-T_L^{\nu_k}-e_kv^{L+1}}{1-\alpha v},
\]
\[
 A_{k,L}'=\frac{\nu_k(1+v)^{\nu_k-1}-(T_L^{\nu_k})'
                -(L+1)e_kv^L+\alpha A_{k,L}}{1-\alpha v},
\]
\[
 A_{k,L}''=\frac{\nu_k(\nu_k-1)(1+v)^{\nu_k-2}-(T_L^{\nu_k})''
          -L(L+1)e_kv^{L-1}+2\alpha A_{k,L}'}{1-\alpha v}.
\]
These identities account for every finite-tail term. If the moment enclosure
is too wide or a local analytic check is inconclusive, the code falls back to
the direct finite recurrence. Neither implementation substitutes the limiting
analytic function for the finite polynomial.

The polynomial corrections are evaluated in their factored, compressed bases,
including the additional 128 blocks of $X$. The tail uses the exact product
rule for $H=\omega q^{1866}(\gamma q^8+E)$. Consequently the evaluated finite jets are
\[
\begin{aligned}
 P&=(z-a)h_L/h_L(a)+D+H,\\
 P'&=[h_L+(z-a)h_L']/h_L(a)+D'+H',\\
 P''&=[2h_L'+(z-a)h_L'']/h_L(a)+D''+H''.
\end{aligned}
\]



\section{Quotient enclosures at the independent base}\label{sec:values}
This section gives the new formulas used by the verifier. Set
\[
 V=P(\astar),\qquad D_\star=P'(\astar),\qquad
 T_*(z)=P_*(z),\qquad g=T_*',
\]
where $V,D_\star$ are directly enclosed by the finite polynomial evaluator.
The initializer proves $|D_\star|>0$. Define
\[
 E_0=|z-a|e_g,\qquad E_1=e_g+|z-a|e_{g'}.
\]
The preceding error bounds give $|P-T_*|\le E_0$ and $|P'-g|\le E_1$.
For a disk $D=\{|z-z_0|\le\rho\}$ avoiding $\astar$, define
\[
 A_*(z)=\frac{T_*(z)-V}{z-\astar},\qquad
 \widetilde Q(z)=\frac{T_*(z)-V}{(z-\astar)D_\star}.
\]
Then
\[
 \widetilde Q'(z)=\frac{g(z)-A_*(z)}{(z-\astar)D_\star},
 \qquad
 \left|\frac{P(z)-V}{(z-\astar)D_\star}-\widetilde Q(z)\right|
 \le\frac{E_0}{|z-\astar||D_\star|}.
\]
Consequently the model-center enclosure has center modulus
$|\widetilde Q(z_0)|$ and error radius
\begin{equation}\label{eq:model-value}
 \rho\sup_D|\widetilde Q'|
 +\sup_D\frac{E_0}{|z-\astar||D_\star|}.
\end{equation}
Subtract this radius for a lower bound and add it for an upper bound.
All suprema and center values are evaluated by outward intervals.

For the finite-center test, use the directly evaluated modulus
\[
 q_0=\left|\frac{P(z_0)-V}{(z_0-\astar)D_\star}\right|.
\]
Since
\[
 \left|P'(z)-\frac{P(z)-V}{z-\astar}\right|
 \le |g(z)-A_*(z)|+E_1+\frac{E_0}{|z-\astar|},
\]
a whole-disk variation bound is
\begin{equation}\label{eq:finite-value}
 R_Q=\rho\sup_D
 \frac{|g-A_*|+E_1+E_0/|z-\astar|}{|z-\astar||D_\star|}.
\end{equation}
The interval $[q_0-R_Q,q_0+R_Q]$ encloses the critical quotient in the disk.
These formulas follow by differentiation and integration along a radius;
they do not assume that the center is an exact critical point.
The interval divisions prove that every accepted disk avoids $\astar$.
Uncertainty in $V$ and $D_\star$ is retained in every operation.



The exhaustive run alone proves
\[
 \BaseLower<\Smal(p,0)<\BaseUpper.
\]
The refinement retests \RefinedCount\ potentially limiting containing disks.
Each contains exactly one root by the exhaustive stage, so its refined and
original quotient intervals can be intersected. All other disks retain their
original bounds. The resulting global enclosure is
\[
 \RefinedLower<\Smal(p,0)<\RefinedUpper.
\]
The minimizing disk has index \MinimizerIndex. For every other critical point,
\[Q(c)>\OtherLower.\]
\UniquenessExplanation
\section{Interval arithmetic and numerical trust assumptions}\label{sec:arithmetic}
The backend encloses basic real operations by rounding the computed
\texttt{long double} endpoint one representable number outward. It runs in
IEEE round-to-nearest mode with binary arithmetic and at least 64 significand
bits. Rational constants are first enclosed with GMP/MPFR conversion.
Nonfinite intervals, zero-containing divisors, and unsupported formats are
rejected. Transcendental endpoints are computed by MPFR at 192 bits with
rounding toward the appropriate infinity~\cite{MPFR}. Compiler flags disable
fast-math reassociation and contraction of multiply-add expressions.

\subsection{Floating-center complex disks}
A complex enclosure is $D(c,r)$, with a floating-point center and an outward
radius. Before center roundoff is added, the exact rules are
\[
 D(c,r)+D(d,s)\subset D(c+d,r+s),\qquad
 D(c,r)D(d,s)\subset D(cd,|c|s+|d|r+rs).
\]
Let $u=2^{-p}$ for significand precision $p$, and let $\tau$ be the least
positive normal number of the hardware type. It dominates the subnormal
rounding floor. If $M_c\ge|c|$ and $M_d\ge|d|$, the implementation adds
\[
 4u(M_c+M_d)+8\tau\quad\hbox{for complex addition},
\]
\[
 8uM_cM_d+32\tau\quad\hbox{for complex multiplication}.
\]
These are conservative upper bounds for the Euclidean error of the computed
center. To verify them, use the scalar round-to-nearest bound
$|\operatorname{fl}(x)-x|\le u|x|+\eta/2$, where $\eta$ is the least positive
subnormal. The two additions have total absolute error at most
$\sqrt2u(M_c+M_d)+\eta$. For a complex product, four products and two additions
give total absolute error at most
$(4u+2u^2)M_cM_d+(3+2u)\eta$. The displayed allowances dominate both for $p\ge64$.
A similar scalar-multiplication bound includes the radius of any interval
scalar. Magnitudes may be replaced by their larger $\ell^1$ bounds; all
radius computations remain outward interval operations.

Horner propagation carries the value and its first two derivatives. For
example, $(Uq+w)''=U''q+2U'q'+Uq''+w''$. Fixed interval coefficients are
converted to enclosing disks; their conversion error is added to the radius.
These rules apply on a complete input box as well as at a center. They avoid
repeated rectangular wrapping but do not change any polynomial coefficient.
The exact-rational regression program performs 64,800 image-containment tests,
including subnormal and extreme-exponent inputs. It supplements this error
argument rather than replacing it.

High powers are evaluated by rotating the input through a quarter turn into
a verified right half-plane and using interval modulus and argument. The
rotation is undone algebraically. All trigonometric argument ranges used by
this certificate are bounded; the endpoint/extremum enumeration in the
interval backend includes every possible extremum on those ranges. Wide
intervals are enclosed by $[-1,1]$ instead.

The calculation relies on these enclosure arguments, the supplied source,
GMP/MPFR, the compiler with the stated flags, and the hardware semantics
specified above.



\section{Complete verification, audits, and reproduction}\label{sec:reproduce}
The exhaustive wrapper requires all sixteen work partitions to pass and
checks that the recorded critical-point indices are exactly
$0,\ldots,29982$, each once. The partitions are index classes modulo sixteen;
they are not geometric sheets used to infer root coverage.
The containing disks are rational data, and their disjointness and the root
counts do not depend on the numerical search.

\file{make_parameters.py} checks the exact coefficient tables and both
base-point coordinates against the native rational input serialization.
\file{run_verification.py} compiles the verifier and runs all partitions.
\file{refine_certificate.py} binds its sharper calculations to the hashes and
complete coverage of the exhaustive stage. \file{audit_certificate.py}
checks degree bounds, data identity, coverage, and the combined intervals.
\file{audit_exact_disks.py} independently checks rational containing-disk
disjointness, accepted-disk containment, and every recorded strict root-count
comparison. \file{audit_normalization.py} additionally checks the new base
parameters and the nonzero derivative enclosures in all partition logs.

The exact-rational audit expands printed endpoints outward. The base summary
widens endpoints by $10^{-15}$, while the refinement and coverage audit use
$(1+|\mathrm{endpoint}|)10^{-18}$. Published endpoints are rounded outward
again. An integrity audit checks recorded calculations; it does not replace
the native analytic enclosure computation.



From the repository directory, run
\begin{lstlisting}
python3 certify.py --jobs 4
\end{lstlisting}
This runs exact input checks, local arithmetic and anchor comparisons,
exhaustive verification, refinement, and the coverage and normalization
checks. It requires Python 3.10 or later, a C++17 compiler, GMP with C++
headers, and MPFR on the supported binary \texttt{long double} platform.
The README gives the build commands and arithmetic requirements.
To check the supplied completed calculation without rerunning its interval
operations, use \file{check_stored_certificate.py}.



\appendix
\section{Exact seed parameters and coefficient tables}
Every terminating decimal in this appendix is exact. The coefficient arrays are in the accompanying JSON specification.
\begingroup\small
\[\Re a=1.0408093642746148137157525156090766412437.\]
\[\Im a=0.3041902165017496705262264286568472099946.\]
\[\Re r=0.9869461588252050688137472794964248376867.\]
\[\Im r=-0.0432448562530681859520443550256337627761.\]
\[\Re b=0.908850253183765306046968000840344747056738618388026628837241.\]
\[\Im b=0.076267057870218869943792292811043863508717811372210553566224.\]
\[\Re\xi_1=0.94945305488086338259789282088083367274676142809337.\]
\[\Im\xi_1=0.01295609450414720552969704525015390110037343638189.\]
\[\Re\xi_2=0.94845012497531595723791829728813409781718653511944.\]
\[\Im\xi_2=0.01327429696527108235832258377560890558082991312789.\]
\endgroup

The normalization base $\astar$, amplitude $\gamma$, and phase parameter are
given in Section~\ref{sec:definition}. The JSON contains all 1,024 coefficients
of $R_0$, four cubic-remainder coefficients, 2,048 coefficients of
$X/(z-a)^2$, and 128 coefficients of $E$. These tables and the finite
recurrence define the polynomial without a monomial expansion.

\begin{thebibliography}{9}
\bibitem{DLMFComplex} NIST Digital Library of Mathematical Functions,
\S1.10(iv), Rouch\'e's theorem, and \S1.11(i), polynomial zeros.
\url{https://dlmf.nist.gov/1.10} and \url{https://dlmf.nist.gov/1.11}.
\bibitem{DLMFBeta} NIST Digital Library of Mathematical Functions,
\S5.12, Euler's beta integral, equation 5.12.1.\ \url{https://dlmf.nist.gov/5.12}.
\bibitem{MPFR} GNU MPFR Manual, rounding modes and correct rounding.
\url{https://www.mpfr.org/mpfr-current/mpfr.html}.
\end{thebibliography}
\end{document}
